\section*{अध्याय \ref{ch:b2:catalytic-cycles} --- कार्बधात्विक उत्प्रेरण: मूल चरण और चक्र}

\begin{solution}{exo:b2:catalytic-cycles:1}
पहले: Ir(I), $d^8$, $8 + 4 \times 2 = 16$ इलेक्ट्रॉन। बाद में: Ir(III), $d^6$, $6 + 6 \times 2 = 18$
इलेक्ट्रॉन, अष्टफलकीय।
\end{solution}

\begin{solution}{exo:b2:catalytic-cycles:2}
पहला \omterm{def:b2:catalytic-cycles:oxidative-addition}{अपचायी विलोपन} है (Pt(IV) से Pt(II), दो मेथिलों से एथेन बनती है)।
दूसरा \omterm{def:b2:catalytic-cycles:ligand-exchange}{लिगैंड विनिमय} है (एक फ़ॉस्फ़ीन का स्थान एथीन लेती है) जिसके बाद
Pd--Ph आबंध में एथीन का \omterm{def:b2:catalytic-cycles:insertion}{प्रवासी निवेशन} होता है।
\end{solution}

\begin{solution}{exo:b2:catalytic-cycles:3}
TON $= 15.0/0.020 = 750$; TOF $= 750/3.0 = \qty{250}{h^{-1}}$।
\end{solution}

\begin{solution}{exo:b2:catalytic-cycles:4}
\begin{align*}
  &\mathrm{PdL_2 + ArBr \to ArPdBrL_2}\\
  &\mathrm{ArPdBrL_2 + Ar'B(OH)_3^- \to ArPdAr'L_2 + Br^- + B(OH)_3}\\
  &\mathrm{ArPdAr'L_2 \to Ar{-}Ar' + PdL_2}
\end{align*}
योग: $\mathrm{ArBr + Ar'B(OH)_3^- \to Ar{-}Ar' + Br^- + B(OH)_3}$; पैलेडियम की हर स्पीशीज़
कट जाती है।
\end{solution}

\begin{solution}{exo:b2:catalytic-cycles:5}
एथिल और आइसोप्रोपिल (इनके $\beta$ हाइड्रोजन हैं)। मेथिल का कोई $\beta$ कार्बन नहीं; नियोपेंटिल के $\beta$
कार्बन पर कोई हाइड्रोजन नहीं; बेन्ज़िल का $\beta$ कार्बन एक \omterm{def:b2:huckel:aromatic}{ऐरोमैटिक} वलय-कार्बन है जिस पर
कोई हाइड्रोजन नहीं।
\end{solution}

\begin{solution}{exo:b2:catalytic-cycles:6}
मेथिल (\textit{E})-3-फ़ेनिलप्रोपीनोएट, \ce{Ph-CH=CH-COOCH3}, फ़ेनिल अंतस्थ
कार्बन पर जुड़ा, trans उत्पाद।
\end{solution}

\begin{solution}{exo:b2:catalytic-cycles:7}\emergencystretch=3em
ब्यूटेनल \ce{CH3CH2CH2CHO} (रैखिक: धातु अंतस्थ कार्बन पर जुड़ती है, हाइड्राइड
भीतरी पर) और 2-मेथिलप्रोपेनल \ce{(CH3)2CHCHO} (शाखित: धातु भीतरी
कार्बन पर)।
\end{solution}

\begin{solution}{exo:b2:catalytic-cycles:8}
\ce{[RhCl(PPh3)3]} के 16 इलेक्ट्रॉन हैं: \ce{H2} जोड़ना (+2) छह
उपसहसंयोजन स्थितियों के साथ 18 देता, जो तीन भारी-भरकम फ़ॉस्फ़ीनों के साथ बहुत भीड़भाड़ वाला है। \ce{[RhCl(PPh3)2]} के 14
हैं और एक खुला स्थल: \omterm{def:b2:catalytic-cycles:oxidative-addition}{ऑक्सीकारी योग} आसानी से 16-इलेक्ट्रॉन संकुल देता है।
\end{solution}

\begin{solution}{exo:b2:catalytic-cycles:9}
बोरोनिक अम्ल दुर्बल नाभिकरागी है; क्षारक बोरॉन पर हाइड्रॉक्साइड (या ऐल्कॉक्साइड) जोड़ता है,
जिससे बोरेट \ce{ArB(OH)3-} बनता है, जिसका ऐरिल समूह अधिक इलेक्ट्रॉन-समृद्ध है और
पैलेडियम पर स्थानांतरित हो जाता है।
\end{solution}

\begin{solution}{exo:b2:catalytic-cycles:10}
20-इलेक्ट्रॉन संकुल तर्कसंगत नहीं है। संकुल को पहले एक लिगैंड खोना चाहिए (14 e), फिर
\ce{H2} जोड़ना (16 e), ऐल्कीन बाँधना (18 e), उसे किसी M--H आबंध में निविष्ट करना (16 e) और
ऐल्केन का विलोपन करना (14 e), इससे पहले कि वह लिगैंड वापस ले या फिर से आरंभ करे।
\end{solution}

\begin{solution}{exo:b2:catalytic-cycles:11}
प्रारंभिक दर $n_{\max}/\tau = \qty{0.25}{mmol/min}$; प्रारंभिक TOF $= 0.25/0.010 =
\qty{25}{min^{-1}} = \qty{1.5e3}{h^{-1}}$। अंतिम TON $= 10.0/0.010 = 1.0 \times 10^3$।
\end{solution}

\begin{solution}{exo:b2:catalytic-cycles:12}
4-ब्रोमोऐनिसोल के साथ फ़ेनिलबोरोनिक अम्ल (या 4-मेथॉक्सीफ़ेनिलबोरोनिक अम्ल के साथ ब्रोमोबेन्ज़ीन),
एक पैलेडियम(0) फ़ॉस्फ़ीन उत्प्रेरक और एक क्षारक। चक्र: ऐरिल ब्रोमाइड का
\omterm{def:b2:catalytic-cycles:oxidative-addition}{ऑक्सीकारी योग}, बोरेट से \omterm{def:b2:catalytic-cycles:transmetalation}{पारधात्वन}, 4-मेथॉक्सीबाइफ़ेनिल का
\omterm{def:b2:catalytic-cycles:oxidative-addition}{अपचायी विलोपन}।
\end{solution}

\begin{solution}{pb:b2:catalytic-cycles:1}
\textbf{1.} Pd(II), $d^8$।
\textbf{2.} Pd(0), $d^{10}$, $10 + 2 \times 2 = 14$ इलेक्ट्रॉन।
\textbf{3.} 14 इलेक्ट्रॉनों के साथ यह दो और युग्म ग्रहण कर सकता है: इसके पास \omterm{def:b2:catalytic-cycles:unsaturated}{रिक्त स्थल} हैं।
\textbf{4.} पैलेडियम(0) फ़ॉस्फ़ीन संकुल वायु से ऑक्सीकृत होते हैं, और फ़ॉस्फ़ीन
फ़ॉस्फ़ीन ऑक्साइडों में ऑक्सीकृत होते हैं।
\textbf{5.} यह पैलेडियम(II) को पैलेडियम(0) में अपचित करती है (स्वयं ऑक्सीकृत होकर) और
पैलेडियम(0) को बाँधकर उसे विलयन में बनाए रखती है।
\textbf{6.} \omterm{def:b2:catalytic-cycles:cycle}{पूर्व-उत्प्रेरक}।
\textbf{7.} \ce{Pd(PPh3)2 + CH3C6H4Br -> [Pd(C6H4CH3)Br(PPh3)2]}: Pd(II), $d^8$, $8 + 4 \times 2 =
16$ इलेक्ट्रॉन।
\textbf{8.} यह उसे बोरेट \ce{PhB(OH)3-} में बदलता है, जो अपना फ़ेनिल समूह स्थानांतरित करता है।
\textbf{9.} \ce{[Pd(Ar)Br(PPh3)2] + PhB(OH)3- -> [Pd(Ar)(Ph)(PPh3)2] + Br- + B(OH)3}।
\textbf{10.} \omterm{def:b2:catalytic-cycles:oxidative-addition}{अपचायी विलोपन} दो cis लिगैंडों को जोड़ता है; trans ऐरिलों को पहले समावयवित होना पड़ता है।
\textbf{11.} \ce{[Pd(Ar)(Ph)(PPh3)2] -> Ar-Ph + Pd(PPh3)2}: पैलेडियम(0), 14 इलेक्ट्रॉन, अगले
चक्कर के लिए तैयार।
\textbf{12.} \ce{CH3C6H4Br + PhB(OH)3- -> CH3C6H4-Ph + Br- + B(OH)3}।
\textbf{13.} 4-ब्रोमोटॉलूईन (\qty{10.0}{mmol}, जबकि बोरोनिक अम्ल के \qty{12.0}{mmol})।
\textbf{14.} $1.51/168.24 = \qty{8.98e-3}{mol}$, अर्थात् \qty{8.98}{mmol}।
\textbf{15.} $8.98/10.0 = 89.8~\%$।
\textbf{16.} $0.050/10.0 = 0.50$~mol~\%।
\textbf{17.} $8.98/0.050 = 180$।
\textbf{18.} $180/4.0 = \qty{45}{h^{-1}}$।
\textbf{19.} एक ही पैलेडियम पर दो फ़ेनिल समूहों का स्थानांतरण (बोरोनिक अम्ल का
समयुग्मन, जिसे ऑक्सीजन के अंश बढ़ावा देते हैं), फिर उनका साथ-साथ विलोपन।
\textbf{20.} ऐरिल समूहों के पास धातु तक पहुँच सकने वाला कोई $\beta$ हाइड्रोजन नहीं है।
\textbf{21.} \omterm{def:b2:catalytic-cycles:oxidative-addition}{ऑक्सीकारी योग}: \ce{C-Cl} आबंध \ce{C-Br} से अधिक प्रबल है।
\textbf{22.} यह अवशेषों में पैलेडियम(0) कणों या लवणों के रूप में समाप्त होता है; यह महँगा और
विषैला है, और पुनः प्राप्त किया जाता है।
\textbf{23.} बोरोनिक अम्ल का एक भाग पार्श्व अभिक्रियाओं में नष्ट होता है (समयुग्मन, वलय से
बोरॉन की हानि); आधिक्य ब्रोमाइड को सीमांत बनाए रखता है।
\textbf{24.} पैलेडियम ने $\boldsymbol{\approx \num{1.8e2}}$ आवर्त पूरे किए।
\end{solution}
